\section{A uniform obstruction to every constant equality offset at radix three}
\label{sec:radix-three-no-go}

The following theorem rules out a concrete class of tempting source edits.
It is not a lower bound on the number of states of a consistency machine,
nor an obstruction to all possible ternary formula encodings.

Let
\[
 p(a,b)=\frac{(a+b)(a+b+1)}2+b,
 \qquad I(A,B)=3p(A,B),
\]
\[
 M(a,b)=3p(a,b)+1,\qquad U(a,B)=3p(B,a)+2,
 \qquad E_e(a,b)=3p(a,b)+e.
\]
Here $e$ is any fixed nonnegative integer. As in the supplied source,
write $N(A)=I(A,1)$ and $\exists_a A=N(U(a,N(A)))$; the refutation target
is the numerical code $1=M(0,0)$. Arrows below mean the numerical
constructor $I$, not an additional proof rule.

\begin{theorem}[No constant equality-tag repair at radix three]
\label{thm:all-offsets}
For every $e\ge0$, the numerical calculus obtained by keeping the supplied
PA proof rules and axiom families, replacing the constructor coefficient
four by three, and replacing the equality tag by $e$, has a finite
derivation of code one. Moreover the derivation is produced uniformly
from $e$ by a primitive-recursive procedure, whose correctness PA proves.
\end{theorem}

The proof uses only the following part of the original calculus:
modus ponens, generalization of closed theorems, the restricted
implicational row
\[
 ((C\to u)\to A)\to((A\to C)\to(B\to C)),\qquad u\le A,
\]
and these heads:
\begin{align*}
 O(d,H)&=(\exists_d U(d,H))\to H,\\
 W(d,h)&=\exists_d E_e(h,d),\\
 L(c,a,b)&=M(c,a)\to(E_e(c,b)\to M(b,a)),\\
 R(c,a,b)&=E_e(c,a)\to(M(b,c)\to M(b,a)),\\
 K_U(d,A,B)&=U(d,A)\to(U(d,I(A,B))\to U(d,B)).
\end{align*}
In the residue-one case with $e\ge4$, it additionally uses membership
vacuity
\[
 V(a,b,c)=M(a,b)\to U(c,M(a,b)),\qquad c\ne a,b,
\]
which is a derived head of the retained guarded gate. Neither diagonal
vacuity nor quantifier commutation is needed by the obstruction.

\paragraph{Common syntactic tools.}
The abstract proof of the restricted implicational basis applies to
$I(A,B)=3p(A,B)$: its only numerical hypotheses are reflexivity of
$\le$ and the bounds $A,B\le I(A,B)$. Thus the fixed $K,S$, composition,
identity and Peirce proof templates remain valid. This follows from the
explicit interface of
\texttt{M6/Logic/RestrictedImplication.lean}; no four-residue decoding
argument enters that proof.

Code zero itself is a theorem: the all-zero instance of the restricted
row has numerical value zero. Also $O$ yields explosion $1\to H$ by a
fixed implicational proof. Under a temporary assumption $1$, $K$ yields
$U(d,N(U(d,H)))\to1$, which is $\exists_d U(d,H)$; applying $O(d,H)$
and discharging the assumption gives $1\to H$. Consequently the usual
classical templates relative to $N$ apply, including extraction of
$V$ from the retained gate $N(I(V,Q))$.

From a closed proof $A\to B$, generalization and $K_U$ yield
$U(d,A)\to U(d,B)$. Applying this construction to the contraposition
$N(B)\to N(A)$ and contraposing again gives existential monotonicity:
\[
 A\to B\quad\Longrightarrow\quad
 (\exists_d A)\to(\exists_d B).
 \tag{EM}
\]
Every use of generalization here is on a closed theorem.

\begin{proof}[Proof of Theorem~\ref{thm:all-offsets}]
We split $e$ into its three possible residues.

\emph{Case 1: $e=3k$, $k\ge0$.}
Let $x=k+1$ and $T=I(x,0)$. Since zero is a theorem, $K(0,x)$ and
modus ponens give a proof of $T$. The bounds above imply $T\ge k+1$,
so the natural number
\[
 n=p(T,1)-k
\]
is well-defined without truncated-subtraction ambiguity. Choose
$(h,d)=p^{-1}(n)$, using the total Cantor inverse. Then
\[
 A:=E_{3k}(h,d)=3(n+k)=3p(T,1)=I(T,1).
\]
There is a fixed $K,S$ proof of $A\to1$: from the theorem $T$, derive
$A\to T$; the identity theorem $A\to A$ is literally
$A\to(T\to1)$; apply $S(A,T,1)$ and modus ponens twice. Generalize
this closed theorem at $d$ to obtain $U(d,N(A))$. But the equality
witness is
\[
 W(d,h)=\exists_d A=N(U(d,N(A))),
\]
so one further modus ponens derives one. This case uses no separation,
vacuity, equality rewriting, or modal head $O$.

\emph{Case 2: $e=3k+2$, $k\ge0$.}
For $k=0$ choose $(h,d)=(1,0)$. For $k\ge1$ choose
$(h,d)=(0,k-1)$. In both cases
\[
 E_{3k+2}(h,d)=U(d,1).
\]
For the second case the required identity is
$p(1,d)-p(0,d)=d+1=k$; the first case is the direct equality
$E_2(1,0)=5=U(0,1)$. Hence the equality-witness head gives
$\exists_d U(d,1)$, and $O(d,1)$ gives its implication to one.
One modus ponens completes this three-line refutation.

\emph{Case 3: $e=3k+1$.}
First let $k=0$. Then $E_1(0,0)=1$. The identity theorem
$N(1)=I(1,1)$, generalized at zero, gives $U(0,N(1))$.
The equality witness $W(0,0)$ is exactly $N(U(0,N(1)))$.
Modus ponens derives one.

Now let $k\ge1$. The pairing identity
\[
 p(a-k,b+k)=p(a,b)+k\qquad(a\ge k)
\]
follows because the sum of the arguments is unchanged. It gives
\[
 E_{3k+1}(2k,0)=M(k,k),\qquad
 E_{3k+1}(3k,k)=M(2k,2k).
\]
The corresponding equality witnesses are
$\exists_0 M(k,k)$ and $\exists_k M(2k,2k)$. Membership vacuity yields
\[
 M(k,k)\to U(0,M(k,k)),\qquad
 M(2k,2k)\to U(k,M(2k,2k)),
\]
with both guards valid because $k\ge1$. Apply (EM), then $O$, to each
witness. This derives both $M(k,k)$ and $M(2k,2k)$. The first of these
is also $E_{3k+1}(2k,0)$. Now the retained rewriting heads give
\[
 M(2k,2k)\to(E_{3k+1}(2k,0)\to M(0,2k))
\]
and
\[
 E_{3k+1}(2k,0)\to(M(0,2k)\to M(0,0)).
\]
Repeated modus ponens therefore derives $M(0,0)=1$.

In every case only fixed proof templates, finite concatenation, numerical
constructor formation, elementary case selection, and Cantor inversion
are used. These operations are primitive recursive. PA proves the displayed
pairing identities, side conditions, and the checking equations of the
fixed templates; induction on their finite proof concatenations gives the
claimed uniform correctness statement.
\end{proof}

\paragraph{Scope of the obstruction.}
The theorem includes the unmodified tag $e=3$, the seemingly natural repair
$e=4$, and every other nonnegative constant offset. It establishes that
merely lowering the coefficient to three and moving the equality tag
cannot preserve the searched consistency theory. A smaller physical table
realizing one of these sources would search an explicitly refutable
calculus. The theorem does not exclude a different nonconstant encoding,
a new axiom system with a proved calibration, or a valid machine with
fewer than 194 states obtained by another construction.
